\section{部署案例研究}

\subsection{案例背景}

某金融客户需要一个 multi-step reasoning agent 来处理合规调查。团队采用 DSPy，将检索、LLM 推理、SQL tool、Verify module 串联起来。部署前的主要风险是 tool 漏洞与 metric drift。

\subsection{控制措施}

\begin{itemize}
    \item 每日 Evidence Snapshot 存入 central repo，并通过 `git tag` 记录 optimizer version
    \item 配置 `Action Timeline` 强制扼制 `fallback path` 早于 production release
    \item 监控 tool success rate，并在每次失败后自动启用 `tool redrive` sequence
    \item 多模型 evaluation：在 A/B 测试中同时跑 two DSPy variants，将 better variant 推向 production
\end{itemize}

\subsection{本章小结}

这个案例展示了 DSPy pipeline 在复杂场景下如何落地：通过 evidence logging、action timeline 与 dual evaluation 把风险控制住。

